	\begin{center}
		\section{基础数学}
	\end{center}

	\begin{multicols*}{2} % 开始双栏，数字 2 表示分两栏

		\subsection{快速幂}


		\begin{minted}{cpp}
// 在 k<=0 的时候返回 1
ll binpow(ll a,ll k,const ll &p=mod){
	ll res=1;
	a%=p;
	while(k>0){
		if((k&1)==1) res=a*res%p;
		a=a*a%p;
		k>>=1;
	}
	return res;
}
		\end{minted}

		\subsection{矩阵快速幂}

		注意，我们这个东西依赖于后面的这个\textcolor{red}{\textbf{自动取模类}}，所以代码还是比较长的。

		\begin{minted}{cpp}
/*
* 2025/11/5 https://www.luogu.com.cn/record/245403611 矩阵快速幂模版题
* 2025/11/5 https://www.luogu.com.cn/record/245411884 P1939 矩阵加速（数列）
* 2026/3/4 去除了过时的宏定义和没有意义的压行
*/
template<class T>
struct Matrix {
	int n, m;
	vector<vector<T> > a; // 1-based

	Matrix() : n(0), m(0) {
	}

	Matrix(int n_, int m_, const T &val = T()) {
		assign(n_, m_, val);
	}

	// 用 1-based 的 vv 构造
	Matrix(const vector<vector<T> > &vv) { load(vv); }

	// vv.size() == n+1, vv[1].size() == m+1
	void load(const vector<vector<T> > &vv) {
		n = (int) vv.size() - 1;
		m = n ? (int) vv[1].size() - 1 : 0;
		a = vv;
	}

	void assign(int n_, int m_, const T &val = T()) {
		n = n_;
		m = m_;
		a.assign(n + 1, vector<T>(m + 1, val));
	}

	// 重载
	vector<T> &operator[](int i) { return a[i]; }
	const vector<T> &operator[](int i) const { return a[i]; }

	// += 负责真正干活
	Matrix &operator+=(const Matrix &rhs) {
		assert(n == rhs.n && m == rhs.m);
		for (int i = 1; i <= n; ++i) {
			for (int j = 1; j <= m; ++j) {
				a[i][j] += rhs[i][j];
			}
		}

		return *this;
	}

	// + 基于 +=
	Matrix operator+(const Matrix &rhs) const {
		Matrix res(*this);
		res += rhs;
		return res;
	}

	Matrix &operator-=(const Matrix &rhs) {
		assert(n == rhs.n && m == rhs.m);
		for (int i = 1; i <= n; ++i) {
			for (int j = 1; j <= m; ++j) {
				a[i][j] -= rhs[i][j];
			}
		}

		return *this;
	}

	Matrix operator-(const Matrix &rhs) const {
		Matrix res(*this);
		res -= rhs;
		return res;
	}

	Matrix &operator*=(const Matrix &rhs) {
		*this = (*this) * rhs;
		return *this;
	}

	Matrix operator*(const Matrix &rhs) const {
		assert(m == rhs.n);
		Matrix res(n, rhs.m, T());
		for (int i = 1; i <= n; ++i) {
			for (int k = 1; k <= m; ++k) {
				if (a[i][k] != T()) {
					for (int j = 1; j <= rhs.m; ++j) {
						res[i][j] += a[i][k] * rhs[k][j];
					}
				}
			}
		}
		return res;
	}

	static Matrix identity(int n) {
		Matrix I(n, n, T());
		for (int i = 1; i <= n; ++i) I[i][i] = T(1);
		return I;
	}

	friend ostream &operator<<(ostream &os, const Matrix &ma) {
		for (int i = 1; i <= ma.n; ++i) {
			for (int j = 1; j <= ma.m; ++j) {
				os << ma[i][j] << " ";
			}
			os << "\n";
		}
		return os;
	}
};

// 矩阵快速幂：必须是方阵
template<class T>
Matrix<T> mat_pow(Matrix<T> base, long long exp) {
	assert(base.n == base.m);
	Matrix<T> res = Matrix<T>::identity(base.n);
	while (exp > 0) {
		if (exp & 1) res = res * base;
		base = base * base;
		exp >>= 1;
	}
	return res;
}
		\end{minted}

		\subsubsection{矩阵快速幂加速递推}
		已知一个数列 $a$，它满足：
		\[
			a_x = \begin{cases}
					  1 & x \in \{1,2,3\} \\ a_{x-1} + a_{x-3} & x \geq 4
			\end{cases}
		\]
		求 $a$ 数列的第 $n$ 项对 $10^9+7$ 取余的值。

		对于递推式 $a_x = a_{x-1} + a_{x-3}$，因为递推涉及到了前 $3$ 项的状态，我们可以构造一个大小为 $3 \times 1$ 的状态列向量 $\begin{bmatrix}
																																   a_x \\ a_{x-1} \\ a_{x-2}
		\end{bmatrix}$。

		根据递推关系，我们可以将其展开为矩阵乘法的形式，寻找转移矩阵 $T$：
		\[
			\begin{aligned}
				a_x     & = 1 \cdot a_{x-1} + 0 \cdot a_{x-2} + 1 \cdot a_{x-3} \\
				a_{x-1} & = 1 \cdot a_{x-1} + 0 \cdot a_{x-2} + 0 \cdot a_{x-3} \\
				a_{x-2} & = 0 \cdot a_{x-1} + 1 \cdot a_{x-2} + 0 \cdot a_{x-3}
			\end{aligned}
		\]
		通过提取系数，可以得出转移矩阵 $T$：
		\[
			\begin{bmatrix}
				a_x \\ a_{x-1} \\ a_{x-2}
			\end{bmatrix} = \begin{bmatrix}
								1 & 0 & 1 \\ 1 & 0 & 0 \\ 0 & 1 & 0
			\end{bmatrix} \begin{bmatrix}
							  a_{x-1} \\ a_{x-2} \\ a_{x-3}
			\end{bmatrix}
		\]
		(这里就是普通的矩阵乘法，按照矩阵乘法规则展开，上式和下式是等价的)
		\\
		初始状态列向量为 $V = \begin{bmatrix}
								  a_3 \\ a_2 \\ a_1
		\end{bmatrix} = \begin{bmatrix}
							1 \\ 1 \\ 1
		\end{bmatrix}$。

		将转移矩阵 $T$ 自乘 $n-1$ 次后，再乘以初始向量 $V$，即可得到最终状态：
		\[
			T^{n-1} V = \begin{bmatrix}
							1 & 0 & 1 \\ 1 & 0 & 0 \\ 0 & 1 & 0
			\end{bmatrix}^{n-1} \begin{bmatrix}
									a_3 \\ a_2 \\ a_1
			\end{bmatrix} = \begin{bmatrix}
								a_{n+2} \\ a_{n+1} \\ a_n
			\end{bmatrix}
		\]
		由此可见，经过 $n-1$ 次状态转移后，底部的元素恰好从 $a_1$ 递推变为了 $a_n$。因此结果矩阵的第 3 行第 1 列即为所求的 $a_n$（在代码中由于加入了空的 \verb|{}| 占位使下标从 1 开始，对应为 \verb|mat[3][1]|）。


		\begin{minted}{cpp}
void Solve() {
	ll n;
	cin >> n;
	vector<vector<modi> > A{
				{},     // 递推矩阵
				{{}, 1, 0, 1},// fn = fn-1 + 0*fn-2 + fn-3
				{{}, 1, 0, 0},
				{{}, 0, 1, 0}
			},
			B = {       // 原向量 v（也可以认为是一个矩阵）
				{},
				{{}, 1}, // a3
				{{}, 1}, // a2
				{{}, 1}  // a1
			};
	Matrix mat(A),matb(B);
	mat = mat_pow(mat, n-1);
	mat*=matb; // 递推矩阵T*原向量 v
	// a1递推n-1次，是这个 an
	cout << mat[3][1] << "\n";
}
		\end{minted}

		对于递推式 $S_{n}=22S_{n-1}-S_{n-2}$，我们可以将其转化为矩阵乘法的形式来加速计算。

		首先构造状态列向量 $\begin{bmatrix}
								S_n \\ S_{n-1}
		\end{bmatrix}$，根据递推关系可以得出转移矩阵 $T$：
		\[
			\begin{bmatrix}
				S_n \\ S_{n-1}
			\end{bmatrix} = \begin{bmatrix}
								22 & -1 \\ 1 & 0
			\end{bmatrix} \begin{bmatrix}
							  S_{n-1} \\ S_{n-2}
			\end{bmatrix}
		\]
		初始状态列向量为 $V = \begin{bmatrix}
								  S_2 \\ S_1
		\end{bmatrix} = \begin{bmatrix}
							482 \\ 22
		\end{bmatrix}$。

		将转移矩阵 $T$ 自乘 $n-1$ 次后，再乘以初始向量 $V$，即可得到最终状态：
		\[
			T^{n-1} V = \begin{bmatrix}
							22 & -1 \\ 1 & 0
			\end{bmatrix}^{n-1} \begin{bmatrix}
									S_2 \\ S_1
			\end{bmatrix} = \begin{bmatrix}
								S_{n+1} \\ S_n
			\end{bmatrix}
		\]
		由此可见，计算结果矩阵的第 2 行第 1 列即为所求的 $S_n$。代码中因为使用了一个空的 \verb|{}| 进行占位（即下标从 1 开始），所以对应的取值为 \verb|T[2][1]|。


		\begin{minted}{cpp}
vector<vector<modi> >
// 构造转移矩阵 T
// [ S_n   ] = [ 22 -1 ] * [ S_{n-1} ]
// [ S_{n-1} ]   [  1  0 ]   [ S_{n-2} ]
// 注意：用 {} 占位使得二维 vector 的下标从 1 开始
tmp={{},
	{{},22,-1},
	{{},1,0},
	};
Matrix T(tmp);

// 将转移矩阵自乘 n-1 次
T=mat_pow(T,n-1);

// 构造初始状态列向量 V
// V = [ S_2 ] = [ 482 ]
//     [ S_1 ]   [  22 ]
vector<vector<modi> > tmp2=
	{{},
	{{},482},
	{{},22},
	};
Matrix V(tmp2);

// 将转移矩阵和初始向量相乘：T^{n-1} * V
T*=V;

// 提取结果 S_n
// 经过 n-1 次转移后，结果列向量变为 [ S_{n+1} ]
//                               [ S_n     ]
// 故所求结果在第 2 行第 1 列（1-indexed）
modi res=T[2][1];
		\end{minted}

		\subsection{自动取模类}
		\begin{minted}{cpp}
// T: 底层存储类型 (uint32_t / uint64_t)
// MOD: 模数
// Wide: 乘法时用的更宽类型
template <class T, T MOD, class Wide = ull>
struct modint {
	T v;
	modint(ll x = 0) {
		ll t = x % (ll)MOD;
		if (t < 0) t += MOD;
		v = (T)t;
	}

	T val() const { return v; }

	modint& operator+=(const modint &o) {
		T x = v + o.v;
		if (x >= MOD) x -= MOD;
		v = x;
		return *this;
	}

	modint& operator-=(const modint &o) {
		T x = v >= o.v ? v - o.v : (T)(v + MOD - o.v);
		v = x;
		return *this;
	}

	modint& operator*=(const modint &o) {
		Wide t = (Wide)v * (Wide)o.v;
		v = (T)(t % MOD);
		return *this;
	}

	// 幂
	static modint powmod(modint a, ll e) {
		modint r = 1;
		while (e) {
			if (e & 1) r *= a;
			a *= a;
			e >>= 1;
		}
		return r;
	}

	// 逆元（假定 MOD 是质数）
	static modint inv(modint a) {
		return powmod(a, (ll)MOD - 2);
	}

	modint& operator/=(const modint &o) {
		return *this *= inv(o);
	}

	// 一元
	modint operator+() const { return *this; }
	modint operator-() const { return modint(0) - *this; }

	// 二元
	friend modint operator+(modint a, const modint &b) { return a += b; }
	friend modint operator-(modint a, const modint &b) { return a -= b; }
	friend modint operator*(modint a, const modint &b) { return a *= b; }
	friend modint operator/(modint a, const modint &b) { return a /= b; }

	friend bool operator==(const modint &a, const modint &b) { return a.v == b.v; }
	friend bool operator!=(const modint &a, const modint &b) { return a.v != b.v; }

	friend ostream& operator<<(ostream &os, const modint &x) {
		return os << x.v;
	}
	friend istream& operator>>(istream &is, modint &x) {
		ll t;
		if (!(is >> t)) return is;
		x = modint(t);
		return is;
	}
};
using modi=modint<uint32_t,mod>;
		\end{minted}

% 我们发现我们自己的那个模板也是线性递推的，因此就不需要这个板子了

% \subsection{预处理逆元（一个数的逆元）}

% 假设要求 $i$ 模 $p$（$p$ 为素数）的逆元，也就是计算 $i^{-1} \pmod p$。

% 首先我们进行带余除法，将 $p$ 分解为包含 $i$ 的形式。令 $q = \lfloor \frac{p}{i} \rfloor$ 为商，$r = p \bmod i$ 为余数，显然有：
% $$
% p = q \times i + r
% $$
% 将等式两边放到模 $p$ 意义下，左边的 $p$ 就变成了 $0$：
% $$
% 0 \equiv q \times i + r \pmod p
% $$
% 因为我们的目标是求出 $i$ 的逆元 $i^{-1}$，我们可以考虑在等式两边同时乘以 $i^{-1} \times r^{-1}$，这样能把 $i$ 和 $r$ 的逆元构造出来：
% $$
% 0 \equiv q \times r^{-1} + i^{-1} \pmod p
% $$
% 移项后，就把未知的 $i^{-1}$ 单独留在一边，用已知的 $r^{-1}$ 表示出来了：
% $$
% i^{-1} \equiv -q \times r^{-1} \pmod p
% $$
% 最后将 $q$ 和 $r$ 的定义代入回来，就得到了完整的递推式：
% $$
% i^{-1} \equiv -\lfloor \frac{p}{i} \rfloor \times (p \bmod i)^{-1} \pmod p
% $$
% 由于余数 $p \bmod i$ 必定严格小于 $i$，因此我们在从小到大递推求解时，$(p \bmod i)^{-1}$ 已经被计算过了，可以直接利用。

% 代码实现中，为了防止 $-q$ 取模产生负数，我们通常将其加上 $p$（即代码中的 \verb|mod - q|）再进行计算。

% 代码参考了 starsilk：

% \begin{minted}{cpp}
% 	inv[1] = 1;
% 	for (int i = 2; i < MAXN - 5; i++) {
% 		long long r = mod % i;
% 		long long q = mod / i;
% 		inv[i] = inv[r] * (mod - q) % mod;
% 	}
% \end{minted}

		\subsection{异或线性基（linear basis）}

		异或线性基用于解决在一个整数集合中，任意选取元素进行异或操作所能产生的数值组合问题。其核心原理基于线性代数中的向量空间，将整数的二进制表示看作位于 $\mathbb{F}_2$ 域（即模 $2$ 加法，等价于异或）上的向量。

		模板中各核心操作的理论基础如下：
		\begin{itemize}
			\item \textbf{插入（\texttt{insert}）}：本质是向极大线性无关组中添加新向量的过程。从高位到低位遍历，若当前数字的第 $i$ 位为 $1$，且该位已存在基底 $d_i$，则将数字异或 $d_i$ 以消去该位；若不存在基底，则将当前数字作为 $d_i$ 存入，插入成功。若最终数字变为 $0$，说明它可由已有基底线性表出。
			\item \textbf{查询最大值（\texttt{query\_max}）}：基于贪心思想。从高位到低位遍历，对于第 $i$ 位的基底 $d_i$，因为它是集合中唯一最高位为 $i$ 的基底，若当前答案异或 $d_i$ 后能变大，则必定选择异或。因为无论后面低位如何选择，都无法弥补当前第 $i$ 位错失的 $1$。
			\item \textbf{查询最小值（\texttt{query\_min}）}：如果在插入时出现过能被线性表出的数字（即能异或出 $0$），那么最小值就是 $0$；否则，由于基底中的元素互相无法表出对方，线性基数组中最小的非零基底就是能异或出的最小值。
			\item \textbf{重构与求第 $k$ 小（\texttt{rebuild} \& \texttt{query\_kth}）}：\texttt{rebuild} 的本质是对线性基矩阵进行高斯消元，化为“简化阶梯型矩阵”。经过重构，对于任意基底 $p_j$，其所有的二进制位 $1$ 在其它基底的对应位上一定都是 $0$。这使得基底之间变成了完全独立的“正交”状态，它们能表出的所有异或和严格按基底的选择与否呈现字典序递增。因此，将排名 $k$ 转化为二进制表示，若第 $j$ 位为 $1$，则异或上第 $j$ 小的非零基底，即可拼凑出第 $k$ 小的异或和。
		\end{itemize}

		\begin{minted}{cpp}
/*
 * 2026-04-06: AC https://vjudge.net/solution/69022982
 * 主要是测试了 rebuild 和 query-kth
 * 2026-04-06: AC https://www.luogu.com.cn/record/272234814
 * 主要测试了 merge 功能和这个 query_max 功能
 * 2026-04-07: AC https://acm.hdu.edu.cn/contest/view-code?cid=1202&rid=3312&from=rank
 * 测试了 query_min_with
 */
struct LinearBasis {
	static const int MAXL = 60;
	vector<long long> d;
	vector<long long> p; // 正交化后的基底，用于求第 k 小
	int cnt; // 线性基内元素的数量（秩）
	bool zero; // 原数组是否能异或出 0

	LinearBasis() {
		d.assign(MAXL + 1, 0);
		cnt = 0;
		zero = false;
	}

	// 向极大线性无关组中添加新向量
	bool insert(long long val) {
		// 从高到低遍历二进制位
		for (int i = MAXL; i >= 0; i--) {
			// 若当前数字的第 i 位为 1
			if (val & (1LL << i)) {
				// 若不存在该位的基底，则作为新基底存入，插入成功
				if (!d[i]) {
					d[i] = val;
					cnt++;
					return true;
				}
				// 若已存在该位基底，则异或消去该位的 1
				val ^= d[i];
			}
		}
		// 若最终数字变为 0，说明它可由已有基底线性表出
		zero = true;
		return false;
	}

	// 查询能异或出的最大值
	long long query_max() {
		long long res = 0;
		// 基于最高位独占特性的贪心，从高位到低位遍历
		for (int i = MAXL; i >= 0; i--) {
			// 若当前答案异或基底 d[i] 后能变大，则必然选择异或
			if ((res ^ d[i]) > res) {
				res ^= d[i];
			}
		}
		return res;
	}

	// 查询能异或出的最小值
	long long query_min() {
		// 若在插入时出现过能被线性表出的数字（即能异或出 0）
		if (zero) return 0;
		// 否则，线性基数组中最小的非零基底就是能异或出的最小值
		for (int i = 0; i <= MAXL; i++) {
			if (d[i]) return d[i];
		}
		return 0;
	}

	// 重构线性基，通过类似高斯消元化为简化阶梯型矩阵
	void rebuild() {
		// 使所有基底互不干扰，达到完全正交状态
		for (int i = MAXL; i >= 0; i--) {
			for (int j = i - 1; j >= 0; j--) {
				// 如果 d[i] 的第 j 位为 1，则用 d[j] 消去这一位
				if (d[i] & (1LL << j)) {
					d[i] ^= d[j];
				}
			}
		}
		p.clear();
		// 将正交化后的非零基底提取出来，方便后续按字典序查询
		for (int i = 0; i <= MAXL; i++) {
			if (d[i]) p.push_back(d[i]);
		}
	}

	// 查询去重后的第 k 小值（要求先调用 rebuild）
	long long query_kth(long long k) {
		// 如果原集合能异或出 0，并且 0 算第一小，先扣除 0 的排名
		if (zero) k--;
		if (k == 0) return 0;
		// 若查询排名超过了能异或出的总个数（2^cnt - 1 个非零值）
		if (k >= (1LL << cnt)) return -1;

		long long res = 0;
		// 将排名 k 转化为二进制表示
		for (int i = 0; i < (int) p.size(); i++) {
			// 若第 i 位为 1，则异或上正交化后的第 i 小的非零基底
			if (k & (1LL << i)) res ^= p[i];
		}
		return res;
	}

	// 合并另一个线性基
	void merge(const LinearBasis &other) {
		for (int i = 0; i <= MAXL; i++) {
			// 将另一个线性基的非零基底全部插入当前线性基
			if (other.d[i]) insert(other.d[i]);
		}
		zero |= other.zero;
	}

	// 给定初值 init，从线性基里选任意子集 XOR，使 init ^ subset_xor 最小
	// 典型场景：1→n 最短 XOR 路径（HDU「月之都」）——先求任意一条路径的异或和
	// init = (1 ^ n) ^ 额外边的 x^y^w，再用环（linearly 入基的 x^y^w）去最小化
	long long query_min_with(long long init) {
		long long res = init;
		// 从高位到低位贪心：d[i] 的最高位就是第 i 位，不会污染更高位；
		// 若 res 的第 i 位为 1，异或 d[i] 能把它消掉，res 严格变小；反之不改
		for (int i = MAXL; i >= 0; i--) {
			if ((res ^ d[i]) < res)
				res ^= d[i];
		}
		return res;
	}
};
		\end{minted}

% // https://www.luogu.com.cn/record/228289095 基础模版
% // https://acm.hdu.edu.cn/contest/view-code?cid=1173&rid=23716
% struct LB{
%     using ll=long long;
%     ll n;vector<ll> b;
%     LB(ll n=62){
%         this->n=n;
%         b.assign(n+5,0);
%     }
%     void insert(ll x){
%         for(int i=n;i>=0;--i){
%             if(x>>i&1){
%                 if(b[i]){
%                     x^=b[i];
%                 }else{
%                     b[i]=x;
%                     return;
%                 }
%             }
%         }
%     }
%     ll findmx(){
%         ll res=0;
%         for(int i=n;i>=0;--i){
%             if((res^b[i]) > res){
%                 res^=b[i];
%             }
%         }
%         return res;
%     }
%     void merge(const LB & t){
%         for(int i=n;i>=0;--i){
%             if(t.b[i]){
%                 insert(t.b[i]);
%             }
%         }
%     }
% };

		\subsection{区间异或线性基（prefix linear basis）}

		处理「区间内任取若干元素异或最大值」的问题。朴素做法是对每个区间重建线性基（$O(n \log V)$ 单次），但不能在线。下面这套\textbf{前缀线性基 + 贪心保留下标靠右元素}的写法可以做到 $O(\log V)$ 单次查询，且天然支持「末尾追加」操作（强制在线场景）。

		模板中各核心操作的理论基础如下：
		\begin{itemize}
			\item \textbf{追加（\texttt{append}）}：对每个前缀 $[1..i]$ 维护一个线性基 $L[i]$，由 $L[i-1]$ 增量得到。插入新元素 $(\text{val}, \text{idx})$ 时，从高位到低位走，对每一个 \text{val} 仍存在的位 $i$：若 $d_i$ 为空，直接落位；否则比较两者下标，将下标更大的那一对存入 $d_i$、$\text{pos}_i$，把下标更小的那一对继续往低位消。这样 $L[r]$ 的每条基底都尽可能用下标靠右的元素来表达。
			\item \textbf{区间查询（\texttt{query\_max}）}：直接取 $L[r]$，从高位到低位贪心，但只接受 $\text{pos}_i \ge l$ 的基底。正确性来源于「靠右优先」的不变量：$L[r]$ 中所有 $\text{pos}_i \ge l$ 的基底，恰好张成 $a[l..r]$ 的异或子空间（每条基底都可写成只用 $a[l..r]$ 内元素的某个异或组合）。
			\item \textbf{空间}：每个 $L[i]$ 是 $\text{BITS}$ 位的基底数组（再加一份 \texttt{pos} 数组），整体 $O(n \cdot \text{BITS})$。值域 $\le 2^{30}$ 时单层 $240$ 字节，$n + m = 10^6$ 约 $240$ MB；若值域更小可压成 \texttt{uint16\_t} 进一步节省。
		\end{itemize}

		\begin{minted}{cpp}
/*
 * 2026-04-29: AC https://vjudge.net/solution/69555764
 * 区间异或线性基模板题：HDU 6579 Operation
 *   op 0 l r：区间 [l, r] 内子集异或最大值（强制在线）
 *   op 1 x  ：在数组末尾追加 x（强制在线）
 */
struct PrefixLinearBasis {
	static constexpr int BITS = 30;
	using T = unsigned;

	struct Layer {
		T   d  [BITS] = {};   // d[i]：最高位为 i 的基底
		int pos[BITS] = {};   // pos[i]：当前 d[i] 对应的元素下标（1-indexed）
	};

	vector<Layer> L;          // L[k]：a[1..k] 的前缀线性基

	PrefixLinearBasis()        { L.emplace_back(); }   // L[0] 为空基
	void reserve(int n)        { L.reserve(n + 1); }
	void clear()               { L.clear(); L.emplace_back(); }
	int  size()  const         { return (int) L.size() - 1; }

	// 末尾追加 val，自动成为 a[size()+1]
	void append(T val) {
		Layer cur = L.back();
		int idx   = (int) L.size();
		for (int i = BITS - 1; i >= 0; --i) {
			if (!((val >> i) & 1)) continue;
			if (!cur.d[i]) {                  // 该位空，落位即可
				cur.d[i] = val;
				cur.pos[i] = idx;
				break;
			}
			// 让 d[i] 始终保留下标更靠右的那一对
			if (cur.pos[i] < idx) {
				swap(cur.d[i], val);
				swap(cur.pos[i], idx);
			}
			val ^= cur.d[i];
		}
		L.push_back(cur);
	}

	// a[l..r] 子集异或最大值（1-indexed，l <= r）
	T query_max(int l, int r) const {
		T ans = 0;
		const Layer &b = L[r];
		for (int i = BITS - 1; i >= 0; --i) {
			if (b.pos[i] < l) continue;       // 该基底用到了 l 之前的元素，跳过
			if ((ans ^ b.d[i]) > ans) ans ^= b.d[i];
		}
		return ans;
	}
};
		\end{minted}

		\subsection{n 进制异或线性基}

		$n$ 进制异或线性基（模 $p$ 域上的线性基）用于处理在 $\mathbb{F}_p$（其中 $p$ 为质数）域上的向量空间问题。与普通的 $\mathbb{F}_2$ 域（二进制异或）类似，在 $\mathbb{F}_p$ 域上，每个数都可以分解为 $p$ 进制的各位数字，构成一个向量。数字间的“$p$ 进制不进位加法”等价于向量的逐元模 $p$ 加法。

		模板中各核心操作的理论基础如下：
		\begin{itemize}
			\item \textbf{转化与消元（\texttt{insert} \& \texttt{query}）}：首先将数字转化为 $p$ 进制向量。利用类似高斯消元法从高到低进行处理，对于每一位，若当前位非零：
			\begin{itemize}
				\item \textbf{插入时}：若该位不存在基底，说明找到了一个新的线性无关向量。为了方便后续消去其他向量的该位，乘以当前最高位元素的乘法逆元，将最高位“归一化”为 $1$ 后存入矩阵（\texttt{mat}），插入成功。若已存在基底，则减去当前位数字与对应基底的乘积，从而消去当前位的非零值。
				\item \textbf{查询时}：逻辑类似。若消元过程中遇到当前位非零但对应位无基底，说明包含无法由当前基表出的独立分量，返回不可表出；若最终整个向量被消为全 $0$，则说明该数字可由基底完全线性表出。
			\end{itemize}
		\end{itemize}

		\begin{minted}{cpp}
/*
 * AC
 * https://acm.hdu.edu.cn/contest/view-code?cid=1199&rid=7820&from=rank
 * n 进制异或线性基（模 p 域）
 */
class modn_XOR_base {
private:
	vector<ll> base; // base[i] 记录最高位为 i 的那个原始数字
	vector<vector<ll> > mat; // mat[i] 记录最高位为 i 的归一化后的 p 进制向量
	ll mod; // 模数 p（必须为质数）

	// 快速幂用于求逆元
	ll power(ll a, ll b) const {
		ll res = 1;
		a %= mod;
		while (b > 0) {
			if (b & 1) res = (res * a) % mod;
			a = (a * a) % mod;
			b >>= 1;
		}
		return res;
	}

	// 费马小定理求乘法逆元（mod 必须为质数）
	ll modInverse(ll n) const {
		return power(n, mod - 2);
	}

public:
	modn_XOR_base(ll Mod) {
		mod = Mod;
		base.assign(64, 0);
		mat.assign(64, vector<ll>(64, 0));
	}

	// 向极大线性无关组中添加新向量（p 进制数字）
	bool insert(ll x) {
		ll orig_x = x;
		vector<ll> v(64, 0);

		// 1. 将数字转化为 p 进制向量（低位到高位）
		for (int i = 0; i < 64; ++i) {
			v[i] = x % mod;
			x /= mod;
		}

		// 2. 从高位到低位遍历，进行高斯消元
		for (int i = 63; i >= 0; --i) {
			if (v[i] == 0) continue; // 当前位为 0，不需要消元

			// 若不存在该位的基底，则作为新基底存入
			if (base[i] == 0) {
				// 求出最高位数字的逆元，用于将最高位归一化为 1
				ll inv = modInverse(v[i]);
				// 整个向量乘以逆元，存入 mat[i]
				for (int j = 0; j <= i; ++j) {
					mat[i][j] = (v[j] * inv) % mod;
				}
				base[i] = orig_x; // 记录原始数字
				return true;
			}

			// 若已存在该位基底，利用对应基底消去当前位的非零值
			ll k = v[i];
			for (int j = 0; j <= i; ++j) {
				// 减去基底对应位的 k 倍（k 刚好是当前最高位的值，因为基底最高位已被归一化为 1）
				v[j] = (v[j] - k * mat[i][j]) % mod;
				if (v[j] < 0) v[j] += mod; // 保证结果为正
			}
		}
		// 若最终向量变为全 0，说明可由已有基底线性表出
		return false;
	}

	// 判断数字 x 能否由当前线性基表出（即 x 是否在张成空间内）
	bool query(ll x) const {
		vector<ll> v(64, 0);

		// 将数字转化为 p 进制向量
		for (int i = 0; i < 64; ++i) {
			v[i] = x % mod;
			x /= mod;
		}

		// 从高位到低位进行消元
		for (int i = 63; i >= 0; --i) {
			if (v[i] == 0) continue;

			// 若当前位不为 0，但基底中对应位为空，说明带有新的独立分量，无法被表出
			if (base[i] == 0) return false;

			// 若存在基底，则继续消元
			ll k = v[i];
			for (int j = 0; j <= i; ++j) {
				v[j] = (v[j] - k * mat[i][j]) % mod;
				if (v[j] < 0) v[j] += mod;
			}
		}

		// 整个向量被成功消为全 0，说明可被线性表出
		return true;
	}
};
		\end{minted}

		\subsection{FFT多项式乘法（非递归）}

		采用非递归版本实现FFT（蝴蝶变换），能显著减小常数并优化内存使用。
		这里通过 \texttt{multiply} 演示了如何运用 \texttt{fft} 对两个多项式（系数数组表示）进行乘法操作。

		\begin{minted}{cpp}
// 理论讲解 https://www.bilibili.com/video/BV1za411F76U/
// 参考（抄的） https://cp-algorithms.com/algebra/fft.html
// 多项式乘法 https://www.luogu.com.cn/record/229525769
using CD = complex<double>;

void fft(vector<CD> &a, bool invert) {
	int n = a.size();
	for (int i = 1, j = 0; i < n; i++) {
		int bit = n >> 1;
		for (; j & bit; bit >>= 1) j ^= bit;
		j ^= bit;
		if (i < j) swap(a[i], a[j]);
	}
	for (int len = 2; len <= n; len <<= 1) {
		double ang = 2 * pi / len * (invert ? -1 : 1);
		CD wlen(cos(ang), sin(ang));
		for (int i = 0; i < n; i += len) {
			CD w(1);
			for (int j = 0; j < len / 2; j++) {
				CD u = a[i + j], v = a[i + j + len / 2] * w;
				a[i + j] = u + v;
				a[i + j + len / 2] = u - v;
				w *= wlen;
			}
		}
	}
	if (invert) {
		for (CD &x : a) x /= n;
	}
}

vector<ll> multiply(vector<ll> &a, vector<ll> &b) {
	vector<CD> fa(a.begin(), a.end()), fb(b.begin(), b.end());
	int n = 1;
	while (n < (int)(a.size() + b.size())) {
		n <<= 1;
	}
	fa.resize(n); fb.resize(n);
	fft(fa, false); fft(fb, false);
	for (int i = 0; i < n; i++) fa[i] *= fb[i];
	fft(fa, true);
	vector<ll> res(n);
	// 避免精度损失
	for (int i = 0; i < n; i++) res[i] = llround(fa[i].real());
	return res;
}
		\end{minted}

		\subsubsection{高精度乘法（基于FFT）}

		https://www.luogu.com.cn/record/229526218

		大整数乘法的核心思想：将大整数的每一位看作多项式的一个系数。例如 $1234$ 可以表示为多项式 $4 + 3x + 2x^2 + x^3$（低位在前）。两个大整数相乘等价于两个多项式相乘，乘完后对结果数组统一处理进位。

		\begin{minted}{cpp}
inline void __() {
	string A, B;
	cin >> A >> B;
	vector<ll> a, b;
	for (int i = (int)A.size() - 1; i >= 0; --i) {
		a.push_back(A[i] - '0');
	}
	for (int i = (int)B.size() - 1; i >= 0; --i) {
		b.push_back(B[i] - '0');
	}
	
	vector<ll> ans = multiply(a, b);
	ll carry = 0;
	for (int i = 0; i < (int)ans.size(); ++i) {
		ans[i] += carry;
		carry = ans[i] / 10;
		ans[i] %= 10;
	}
	while (!ans.empty() && ans.back() == 0) {
		ans.pop_back();
	}
	if (ans.empty()) {
		cout << 0;
		return;
	}
	for (int i = (int)ans.size() - 1; i >= 0; --i) {
		cout << ans[i];
	}
}
		\end{minted}

		\subsubsection{FFT 的使用方法与常见应用}

		\texttt{multiply(a, b)} 的本质是：给定多项式 $A(x) = \sum a_i x^i$ 和 $B(x) = \sum b_i x^i$，返回乘积多项式 $C(x) = A(x) \cdot B(x)$ 的系数数组，其中 $c_k = \sum_{i+j=k} a_i \cdot b_j$。因此，\textbf{只要问题能转化为这种卷积形式，就能用 FFT 加速到 $O(n \log n)$}。

		\textbf{1. 加法卷积计数}

		最典型的场景：有两组数，问它们各取一个数相加能得到每个和的方案数。

		例如掷两颗骰子，求每种点数之和的方案数。将骰子 A 的各面出现次数存入 \texttt{a[]}（\texttt{a[k]} 表示点数 $k$ 出现的次数），骰子 B 同理存入 \texttt{b[]}，则 \texttt{multiply(a, b)} 的结果 \texttt{c[k]} 就是和为 $k$ 的方案数。

		\begin{minted}{cpp}
// 例：两个数组各取一个数，统计每种和的方案数
// a[i] = 数 i 在第一组中出现的次数
// b[j] = 数 j 在第二组中出现的次数
// c[k] = 和为 k 的方案数 = sum(a[i]*b[k-i])
vector<ll> c = multiply(a, b);
		\end{minted}

		\textbf{2. 字符串匹配（含通配符）}

		给定文本串 $T$（长度 $n$）和模式串 $P$（长度 $m$），对于 $T$ 的每个起始位置 $t$，定义失配度：
		$$D(t) = \sum_{j=0}^{m-1} (T_{t+j} - P_j)^2 \cdot w_j$$
		其中 $w_j$ 是通配符权重（通配符处 $w_j = 0$）。$D(t) = 0$ 表示位置 $t$ 完全匹配。

		将 $P$ 翻转为 $P'_j = P_{m-1-j}$，展开后各项都可化为卷积形式，分别用 FFT 计算后合并。

		\begin{minted}{cpp}
// 无通配符的简单情形：检查 T 和 P 是否在每个位置匹配
// 构造 a[i] = T[i], b[j] = P[m-1-j]（翻转 P）
// 则 c[t+m-1] = sum(T[t+j] * P[m-1-j]) 即为位置 t 的内积
// 更一般地，含通配符时需要构造多组卷积分别计算
vector<ll> a(n), b(m);
for (int i = 0; i < n; i++) a[i] = T[i] - 'a' + 1;
for (int j = 0; j < m; j++) b[j] = P[m - 1 - j] - 'a' + 1;
vector<ll> c = multiply(a, b);
// c[t + m - 1] 即为位置 t 的匹配内积
		\end{minted}

		\textbf{3. 高精度乘法}

		如上一节所述，将大整数每一位作为多项式系数，调用 \texttt{multiply} 后统一处理进位。

		\textbf{4. 多项式求值与插值的中间步骤}

		分治 FFT 可以加速多项式多点求值、多点插值等操作，核心都是递归地调用 \texttt{multiply} 合并子问题的结果多项式。

		\subsection{MEX (Minimum Excluded Value)}

		$\operatorname{MEX}$（Minimum Excluded Value）通常指一个非负整数集合中未出现的最小非负整数。
%
%		\subsubsection{单步替换 $=$ 一次销毁 $+$ 一次创生：双台账并行结算原则}
%
%		许多 $\operatorname{MEX}$ 类题目的操作形如"将某位置上的元素替换为某区间的 $\operatorname{MEX}$"。\textbf{不要}把注意力停留在"MEX 是什么"这一层——真正可迁移的抽象是：这个操作是一个\textbf{原子的 $(-x,\, +y)$ 对}，每一次必然严格销毁一个旧元素、创生一个新元素，销毁与创生 \textbf{一次打包、不可拆分}。一旦目标以"最终数组需呈现某种配置"给出，需求便自然分裂为两本\textbf{相互独立}的台账：
%
%		\begin{description}
%			\item[\textbf{销毁台账 $D$}] 目标态下\textbf{不应存在}的元素总数（例如欲使全局 $\operatorname{MEX} = K$，则原数组中所有值为 $K$ 的位置都要改掉）。
%			\item[\textbf{创生台账 $C$}] 目标态下\textbf{必须出现但当前缺失}的元素种数（例如 $0\dots K-1$ 中尚未出现的种类数）。
%		\end{description}
%
%		\textbf{并行结算原则}：若每次操作必定同时向两本台账各记一笔，且"销毁谁 / 创生谁"可以\textbf{独立}挑选（自由度足够），则最少操作次数为
%		\[
%			\operatorname{ans} \;=\; \max(D,\, C)
%		\]
%		而非 $D + C$。因为前 $\min(D,C)$ 笔可被同一次操作\textbf{同时冲抵}（一石二鸟），超出的 $|D-C|$ 笔则由另一侧"白嫖空位"吸收——多出来的销毁可随便作用到无关位置上，多出来的创生可随便写到无关位置上。下界方面，$D$ 与 $C$ 都是\textbf{各自独立}的必要下界（销毁每次至多 $-1$、创生每次至多 $+1$），所以 $\max(D, C)$ 天然是下界；上述"一石二鸟"的构造说明它也可达。
%
%		\begin{center}
%			\begin{tikzpicture}
%				[
%				Dcell/.style={draw, minimum size=6mm, inner sep=0pt, font=\scriptsize\sffamily, fill=red!15},
%				Ccell/.style={draw, minimum size=6mm, inner sep=0pt, font=\scriptsize\sffamily, fill=teal!18},
%				blank/.style={draw=none, fill=none, minimum size=6mm, inner sep=0pt},
%				pair/.style={thick, violet, <->, >=Stealth, shorten <=1pt, shorten >=1pt},
%				]
%				\matrix (m) [
%					matrix of nodes,
%					column sep=3pt,
%					row sep=10mm,
%					row 1/.style={nodes={Dcell}},
%					row 2/.style={nodes={Ccell}},
%				] {
%					$D_1$ & $D_2$ & $D_3$ & $D_4$ & $D_5$ \\
%					$C_1$ & $C_2$ & $C_3$ & |[blank]| & |[blank]| \\
%				};
%				\foreach \i in {1,2,3} \draw[pair] (m-1-\i) -- (m-2-\i);
%
%				\node[left=4pt of m-1-1, font=\scriptsize\sffamily, text=red!70] {销毁 $D$};
%				\node[left=4pt of m-2-1, font=\scriptsize\sffamily, text=teal] {创生 $C$};
%
%				\draw[decorate, decoration={brace, amplitude=3pt}, thin, violet]
%				(m-1-1.north west) -- (m-1-3.north east)
%				node[midway, above=4pt, font=\scriptsize\sffamily, text=violet] {$\min(D,C)$};
%				\draw[decorate, decoration={brace, amplitude=3pt}, thin, red!60]
%				(m-1-4.north west) -- (m-1-5.north east)
%				node[midway, above=4pt, font=\scriptsize\sffamily, text=red!70] {$|D-C|$};
%			\end{tikzpicture}
%
%			\vspace{2pt}
%			{\footnotesize\sffamily\color{gray} 图注：以 $D=5$、$C=3$ 为例。左侧 $\min(D,C)=3$ 笔每次同时从两本台账各划去一项（一石二鸟，紫色双向箭头）；右侧 $|D-C|=2$ 笔只在 $D$ 侧有效，$C$ 侧白嫖无关位置。总操作数 $=\max(D,C)=5$。}
%		\end{center}
%
%		\textbf{自由度的来源——为何 MEX 几乎总满足前提}：
%
%		\begin{itemize}
%			\item \textbf{销毁侧自由（任选位置）}：想抹去哪一位，就选一个覆盖该位置的区间即可，作用下标完全由我们挑。
%			\item \textbf{创生侧自由（任选 MEX 值）}：一次操作能造出任何"当前缺的目标值 $v$"，严格论证见下方专段。
%			\item \textbf{下界的自举（从 0 开始单调扩张）}：任何非零单点的 $\operatorname{MEX}$ 均为 $0$，所以第一步总能凭空得到 $0$；有了 $0\dots k-1$ 就能造 $k$。可达 $\operatorname{MEX}$ 集 $\{0,1,2,\dots\}$ 如搭积木般单调扩张，\textbf{天然与"创生台账需按 $0,1,\dots,K-1$ 顺序补齐"的节奏对齐}，不会出现"想创 $v$ 却因前置未备齐而被迫多花一步"。
%		\end{itemize}
%
%		\textbf{为什么"创生侧自由"真的自由——两条看似严苛的约束都是目标状态白送的}：
%
%		要让一次操作的区间 $\operatorname{MEX}$ 等于指定的 $v$，表面上需要所选子区间\textbf{同时}满足两条约束：
%
%		\begin{description}
%			\item[约束 (i)] 子区间\textbf{含 $\{0,1,\dots,v-1\}$} 的全部 $v$ 个值；
%			\item[约束 (ii)] 子区间\textbf{不含 $v$}。
%		\end{description}
%
%		泛泛看这两条都不容易凑齐（尤其 (ii)，要在一段连续区间里精确避开某个值）。\textbf{但"创生台账待办"这个语境让它们双双白送}：
%
%		\begin{itemize}
%			\item \textbf{(ii) 由目标状态白送}：我们之所以把 $v$ 列入"待创生"，\textbf{恰恰是因为 $v$ 此刻整个数组都不存在}（否则就不算缺）。$v$ 在\textbf{全局}都找不到，任何子区间当然也不会含它——(ii) 对我们\textbf{形同无物}，连"要绕开 $v$ 的位置"都不用考虑。
%			\item \textbf{(i) 由台账递进节奏白送}：由下界自举性质，我们是按 $0,1,2,\dots$ 的顺序\textbf{单调递增}地补齐缺值的，所以轮到要创生 $v$ 时，$\{0,\dots,v-1\}$ \textbf{已经全部在数组中某些位置存在}（或原本就有、或前几步操作刚造出来）。既然这 $v$ 个值都躺在数组里，只要子区间\textbf{大到把它们的位置全部包住}，(i) 就成立——\textbf{"干脆取整个数组"永远是合法保底}，不用费心算具体该取 $[l,r]$ 中的哪两端。
%		\end{itemize}
%
%		\textbf{精髓一句}：\textbf{"想创生的值"与"MEX 要排除的值"恰好指的是同一个 $v$}。目标状态自己送"不含 $v$"，台账节奏自己送"含 $0\dots v-1$"——\textbf{两条约束不是被 MEX 精巧凑出来的，而是被目标与节奏顺手送上门的}。MEX 只是把这个既成事实翻译成了一次合法的产出动作。再叠加销毁侧"任选位置"的自由——把想销毁的废点位置一并纳入这个够大的子区间——\textbf{销毁与创生便在同一次操作内一并兑现}，双台账并行结算原则在 MEX 问题上就能\textbf{零损耗}地生效。
%
%		\begin{center}
%			\begin{tikzpicture}[pairarc/.style={thick, line cap=round, ->, >=Stealth}]
%			\matrix (m) [
%				matrix of nodes,
%				column sep=4pt,
%				row sep=16pt,
%				row 1/.style={nodes={font=\sffamily, inner sep=2pt}},
%				row 2/.style={nodes={font=\sffamily, text=violet, inner sep=2pt}},
%			] {
%				$\dots$ & 0 & 1 & $\dots$ & $X\!-\!1$ & $K$ & $\dots$ \\
%				& & & & & $X$ & \\
%			};
%
%			\begin{scope}[on background layer]
%				\fill[teal!15, rounded corners] (m-1-2.north west) rectangle (m-1-6.south east);
%			\end{scope}
%
%			\draw[decorate, decoration={brace, mirror, amplitude=4pt}, thick, teal]
%			(m-1-2.south west) -- (m-1-5.south east)
%			node[midway, below=6pt, font=\scriptsize\sffamily, text=teal] {$\operatorname{MEX}=X$};
%
%			\draw[pairarc, violet, dashed] (m-1-6.south) -- (m-2-6.north)
%			node[midway, right=1pt, font=\scriptsize\sffamily, text=violet] {替换};
%			\end{tikzpicture}
%
%			\vspace{2pt}
%			{\footnotesize\sffamily\color{gray} 图注：框选包含 $0 \dots X-1$ 且不含 $X$ 的区间，其 $\operatorname{MEX}$ 恰为 $X$。一次替换既消灭废点 $K$，又创造了 $X$。}
%		\end{center}
%
%		\textbf{实战范例}（\href{https://www.codechef.com/START235B/problems/P5235}{CodeChef MEX Spectrum}）：每次可选一个区间，将其中的一个元素替换为该区间的 $\operatorname{MEX}$，求使全局 $\operatorname{MEX} = K$ 的最少操作次数。直接套用上面的两本台账：$D = \operatorname{cnt}[K]$（原数组中所有等于 $K$ 的位置都得抹去），$C = K - \operatorname{havCnt}$（其中 $\operatorname{havCnt}$ 为 $0\dots K-1$ 中出现过的种数）。由并行结算原则，答案即 $\max(\operatorname{cnt}[K],\; K - \operatorname{havCnt})$，$O(N)$ 单次扫描即可对所有 $K$ 同时输出。\textbf{赛时能"猜"出这个结论，本质就是下意识捕捉到了 $D$、$C$ 两条下界；把它上升为上述"双台账并行结算原则"后，再遇到类似的单步替换型构造题（不限于 MEX），就能直接套路化地判断何时取 $\max$、何时需要加权和，不必每题重猜。}

		\subsection{中位数 (Median)}

		本节约定 \textbf{长度奇数} $N$ 的数组，中位数即排序后的第 $\lceil N/2 \rceil$ 个元素（也等于第 $(N+1)/2$ 个）。中位数题的"特征气味"：题面里提到 \textbf{排序后取中间值} / \textbf{划分若干段、每段中位数相同} / \textbf{操作不改变中位数} 之类。

		\subsubsection{中位数的双向判定不等式}

		"$m$ 是数组 $A$ 的中位数"的等价条件是\textbf{两条对称的下界同时成立}：

		\[
			\textcolor{blue}{\bm{
				\#\{a \le m\} \ge \lceil N/2 \rceil \quad \text{且} \quad \#\{a \ge m\} \ge \lceil N/2 \rceil
			}}
		\]

		\textbf{为什么是"$\le$"和"$\ge$"，而不是"$<$"和"$>$"}：要让 $m$ 这个具体的值\textbf{出现在中间位置}，必须\textbf{同时}有 $\ge \lceil N/2 \rceil$ 个不超过它的、$\ge \lceil N/2 \rceil$ 个不少于它的——两条等号都包含 $m$ 自己，正是它落在中间的根据。如果换成严格不等式，只能限定中位数 $\ge m$ 或 $\le m$ 单边，得不到 $=m$。

		\textbf{推论：包含性自动达成}。两条相加，等号情形给出
		\[
			\#\{a \le m\} + \#\{a \ge m\} \;=\; N + \#\{a = m\},
		\]
		右边减 $N$ 等于左边减 $N \ge 2\lceil N/2\rceil - N = 1$（奇 $N$），所以 $\#\{a=m\} \ge 1$。\textbf{即"$m$ 在数组里出现至少一次"是这两条下界的\textbf{推论}，不是独立条件}。判区间中位数等于 $m$ 时\textbf{不必单独判 $m$ 在不在}，套这两条不等式即可——这是常踩的坑。

		\subsubsection{$\pm1$ 编码 $+$ 前缀和：$O(1)$ 区间中位数判定}

		要 $O(1)$ 判 \enquote{$\operatorname{med}(a[l..r]) = m$}（区间长度奇），定义两套 $\pm 1$ 编码：

		\[
			b_1[k] = \begin{cases}
						 +1, & a_k \le m \\ -1, & a_k > m
			\end{cases},
			\qquad
			b_2[k] = \begin{cases}
						 +1, & a_k \ge m \\ -1, & a_k < m
			\end{cases}
		\]

		分别取前缀和 $B_1, B_2$。对于奇长 $\text{len} = r - l + 1$：

		\begin{align*}
			B_1[r] - B_1[l-1] \;&=\; \#\{a \le m\} - \#\{a > m\} \\
			\;&=\; 2\,\#\{a \le m\} - \text{len}.
		\end{align*}

		所以 $\#\{a \le m\} \ge \lceil \text{len}/2 \rceil = (\text{len}+1)/2$ 当且仅当 $B_1$ 差值 $\ge 1$。$B_2$ 同理。综合：

		\[
			\textcolor{blue}{
				\begin{aligned}
					\operatorname{med}(a[l..r]) = m \;\iff\; & B_1[r] - B_1[l-1] \ge 1 \\
					\text{且}\;\,& B_2[r] - B_2[l-1] \ge 1
				\end{aligned}
			}
		\]

		预处理 $O(N)$，单次判 $O(1)$。这个 $\pm 1$ + 双前缀和的组合是中位数题的\textbf{瑞士军刀}：除了等于判定，还能拓展到"区间中位数 $\ge m$ 个数"（只要一条 $B_1$ 差 $\ge 1$）和"区间中位数 $\le m$ 个数"（只要一条 $B_2$ 差 $\ge 1$）。

		\subsubsection{多段奇长度划分公共中位数 $=$ 全局中位数}

		\textbf{命题}：长度奇 $N$ 的数组 $a$ 被划分成 $p$ 段奇长度子数组，且\textbf{每段中位数都等于公共值 $m$}，则 $m$ 必等于 $a$ 的整体中位数。

		\textbf{一行证明}：每段 $S_i$ 长 $\ell_i$（奇）、中位数 $m$，由判定不等式 $\#_{S_i}(\le m) \ge (\ell_i+1)/2$。对所有 $p$ 段相加：
		\[
			\#_a(\le m) \;\ge\; \sum_{i=1}^{p} \frac{\ell_i+1}{2} \;=\; \frac{N + p}{2} \;\ge\; \lceil N/2 \rceil.
		\]
		$\#_a(\ge m)$ 同理。两条都成立 $\Rightarrow$ $m$ 是 $a$ 的中位数；奇 $N$ 中位数唯一，所以 $m$ 即整体中位数。\textbf{推论}：处理这类划分题时\textbf{不必枚举公共中位数候选}——直接 $O(N \log N)$ 排序取出 $a_{\lceil N/2\rceil}$ 当 $m$ 即可，问题立刻退化成\enquote{已知 $m$，最大化奇长度段数}。

		\textbf{实战范例}（\href{https://vjudge.net/problem/CodeForces-2222C}{CF 1094C \emph{Median Partition}}）：把奇长 $N$ 数组划分成最多奇长度子数组、每段中位数相同。先求出全局中位数 $m$（上一条命题），再做 $O(N^2)$ DP：$f[i] = \max\{f[j] + 1 : 0 \le j < i,\; (i-j)\text{ 奇},\; \operatorname{med}(a[j{+}1..i]) = m\}$，区间中位数判定走 $\pm 1$ 双前缀和（上一条工具）。两件套合用，$\sum N^2 \le 5000^2$ 直接通过。\textbf{这道题的"难点"完全集中在"公共中位数 $=$ 全局中位数"和"$\pm 1$ 双前缀和"这两个性质上——把性质内化进板子后，类似题型可以"看到中位数 $=$ 公共值"就直接套出 DP 框架，不再需要赛时重新发明判定。}

		\subsubsection{$L_1$ 最优化是中位数（对比 $L_2$ 平方和最优化是平均数）}

		\textbf{命题}：对实数集 $\{a_1, \ldots, a_n\}$，函数 $f(x) = \sum_{i=1}^{n} |a_i - x|$ 在 $x$ 取 $\{a_i\}$ 的\textbf{中位数}时达到最小值。

		\textbf{直觉证明}：把 $x$ 从极小往极大扫，$f$ 的导数（除若干断点外存在）是 $\#\{a_i < x\} - \#\{a_i > x\}$。导数 $\le 0 \Leftrightarrow \#(a_i > x) \ge \#(a_i < x)$ 即 $x$ 在中位数下方；导数 $\ge 0$ 同理。两侧夹逼出"中位数处导数变号"。$n$ 偶时整段 $[a_{n/2}, a_{n/2+1}]$ 都是最小值点。

		\textbf{怎么用}：题面看到\enquote{选一个 $x$ 让 $\sum|a_i-x|$ 最小} / \enquote{把所有人移到同一个位置，移动总距离最小} / \enquote{涂色 / 校准 / 对齐成同一值的最小代价}\textbf{都是中位数}。$O(N \log N)$ 排序取中位即可，不要傻乎乎地三分 / 二分。

		\textbf{对比 $L_2$（平方和）}：$\sum (a_i - x)^2$ 最小值点是\textbf{平均值}。中位数 vs 平均的根本差异就在 $L_1$ vs $L_2$。

		\textbf{实战范例}：\href{https://vjudge.net/problem/NBUT-1569}{NBUT 1569 \emph{|HELLO - WORLD|}}（最小化 $\sum |a_i - x|$ 模板题）、\href{https://vjudge.net/problem/POJ-3579}{POJ 3579 \emph{Median}}（求所有数对绝对差的中位数，二分答案 + 双指针）、\href{https://vjudge.net/problem/Gym-104386B}{Gym 104386B \emph{Random Array}}（同模型变体）。

		\subsubsection{流式 / 增量中位数：对顶堆}

		要支持\textbf{在线 insert + 实时查询当前中位数}的场景，标准做法是\textbf{大根堆 + 小根堆}对顶维护：

		\begin{itemize}
			\item \textbf{下半堆 $L$}：大根堆，存较小的 $\lceil n/2 \rceil$ 个数；堆顶就是中位数（$n$ 奇时）或下中位数（$n$ 偶时）。
			\item \textbf{上半堆 $R$}：小根堆，存较大的 $\lfloor n/2 \rfloor$ 个数；堆顶是上中位数。
			\item \textbf{不变量}：$|L| - |R| \in \{0, 1\}$ 且 $\max L \le \min R$。
		\end{itemize}

		\textbf{insert 流程}：先把新元素和 $\max L$ 比较丢进对应堆；若两堆 size 失衡（差 $\ge 2$）就把多的那堆堆顶搬到另一堆。每次 $O(\log n)$。\textbf{查询中位数}：直接 $\max L$（$n$ 奇）或 $\frac{\max L + \min R}{2}$（$n$ 偶），$O(1)$。

		\textbf{进阶：滑动窗口中位数} 也是同样的对顶堆 + 延迟删除（lazy delete），或者用 \texttt{multiset} + 迭代器维护。复杂度 $O(N \log K)$。

		\textbf{实战范例}：\href{https://vjudge.net/problem/Baekjoon-1655}{Baekjoon 1655 \emph{说出中间值}}（流式中位数模板）、\href{https://vjudge.net/problem/HackerRank-find-the-running-median}{HackerRank \emph{Find the Running Median}}、\href{https://vjudge.net/problem/FZU-1702}{FZU 1702 \emph{query}}。

		\subsubsection{区间第 $k$ 大（中位数是 $k = \lceil \text{len}/2 \rceil$ 的特例）}

		\textbf{静态}区间中位数 / 第 $k$ 大查询，多次询问、不带修改：

		\begin{description}
			\item[\href{https://oi-wiki.org/ds/persistent-seg/}{主席树（可持久化权值线段树）}] 预处理 $O(N \log N)$、单次 $O(\log N)$，最常用。
			\item[整体二分] 二分答案 $\text{mid}$，把所有 $\le \text{mid}$ 的数加进 BIT 算每个询问区间内 $\le \text{mid}$ 的个数，按是否 $\ge k$ 分流。$O((N + Q) \log^2 N)$ 但常数小。
			\item[莫队 + 值域分块] $O((N + Q) \sqrt N)$，写起来短，对询问规模友好。
		\end{description}

		\textbf{动态}（带修改的区间中位数）则要带修主席树 / 树套树，复杂度 $O(\log^2 N)$。

		\textbf{实战范例}：\href{https://vjudge.net/problem/SPOJ-MKTHNUM}{SPOJ MKTHNUM \emph{K-th Number}}（主席树模板）、\href{https://vjudge.net/problem/Baekjoon-7469}{Baekjoon 7469 \emph{第 K 个数}}（同款）。

%		\subsubsection{中位数题的练习题谱}
%
%		按"用到的核心工具"分类，下表把上面三条性质 / 工具映射到原题，赛时识别题型直接套：
%
%		\begin{description}
%			\item[$\pm 1$ 编码 + 前缀和（数子数组）] 把 \enquote{$=m$ 判定}的不等式延伸成 \enquote{$\ge m$ 计数} / \enquote{$=m$ 计数}：把每个元素映射成 $+1$（满足条件）或 $-1$，则\textbf{中位数 $\ge m$} 等价于前缀和差 $\ge 1$、\textbf{中位数 $= m$} 用容斥 $f(\ge m) - f(\ge m+1)$ 即可。计数 \enquote{前缀和差 $\ge 1$ 的对数} 走 BIT / 归并排序 / 双指针均可。代表题：
%			\begin{itemize}[leftmargin=1.2em, topsep=2pt, itemsep=1pt]
%				\item \href{https://vjudge.net/problem/CodeForces-1005E1}{CF 1005E1 \emph{Median on Segments (Permutations Edition)}}：排列里数 \enquote{中位数 $= m$} 的子数组个数，$\pm 1$ 编码模板题。
%				\item \href{https://vjudge.net/problem/CodeForces-1005E2}{CF 1005E2 \emph{Median on Segments (General)}}：上一题的可重复元素版本，$\pm 1$ 编码加一点点小心处理 ties。
%				\item \href{https://vjudge.net/problem/洛谷-P3031}{洛谷 P3031 \emph{Above the Median G}} / \href{https://vjudge.net/problem/HDU-4908}{HDU 4908 \emph{BestCoder Sequence}}：中位数 $\ge m$ 的子数组数，BIT 数对题。
%			\end{itemize}
%
%			\item[公共中位数 $=$ 全局中位数（划分型）] 一旦题面是 \enquote{划分若干奇长段、每段中位数都相同} 这个范式，先排序锁定 $m$，再 DP / 贪心。代表题：
%			\begin{itemize}[leftmargin=1.2em, topsep=2pt, itemsep=1pt]
%				\item \href{https://vjudge.net/problem/CodeForces-2222C}{CF 1094C \emph{Median Partition}}（本节实战范例）：奇长 $N$，最大化段数。$O(N^2)$ DP。
%				\item \href{https://vjudge.net/problem/Gym-105242L}{Gym 105242L \emph{Median of the Array}}：CF 1094C 的"二段简化版"——只问能否切成两段同中位数。同样先求全局中位数，再判存在性。入门顺手题。
%			\end{itemize}
%
%			\item[中位数判定不等式（双向夹逼）] \enquote{$\#(\le m) \ge \lceil N/2 \rceil$ 且 $\#(\ge m) \ge \lceil N/2 \rceil$} 这个判定本身常作为构造 / 计数题的"约束翻译"——只要题面里出现 \enquote{排序后中间值是某个特定数}，立刻翻译成两条计数下界。代表题：
%			\begin{itemize}[leftmargin=1.2em, topsep=2pt, itemsep=1pt]
%				\item \href{https://vjudge.net/problem/qoj-13877}{QOJ 13877 \emph{Arcane Secret}}：$N$ 数等分成 $K$ 组、取每组最大、再求这些最大的中位数；问哪些原数能成为这种"max 中位数"。把 \enquote{$x$ 是中位数} 拆成两条计数下界，再分类讨论 $x$ 在哪一组当 max。
%				\item \href{https://vjudge.net/problem/SPOJ-EQUI}{SPOJ EQUI \emph{Equilibrium}}：往集合里加一个数，让 mean $=$ median；中位数判定 + 简单代数。
%			\end{itemize}
%		\end{description}
%
%		\textbf{识别套路}：题面提到 \enquote{中位数} 时第一时间问自己——"我要的是\textbf{判定 $=m$}（套 $\pm 1$ 双前缀和）、\textbf{计数 $\ge m$ 的子数组}（$\pm 1$ 单前缀和 + BIT）、\textbf{划分若干段同中位数}（先求全局中位数）、还是\textbf{把 $x$ 是中位数转成两条计数约束}（双向夹逼）"。这四条工具基本覆盖中位数题的 80\%，剩下 20\% 才是真正需要单独想的（如可持久化线段树查区间第 $k$ 大、对顶堆维护流式中位数等）。

		\subsection{常见结论（conclusion）}

		\subsubsection{约数个数定理}

		设正整数 $n$ 的质因数分解为：
		$$n = p_1^{e_1} \cdot p_2^{e_2} \cdots p_k^{e_k}$$

		则 $n$ 的约数个数 $d(n)$ 为：
		$$d(n) = (e_1 + 1)(e_2 + 1) \cdots (e_k + 1)$$

		例题：\href{https://atcoder.jp/contests/abc383/tasks/abc383_d}{ABC-383-D 9 Divisors}

		要让 $n$ 恰好有九个约数，由于 $9 = 9 \times 1 = 3 \times 3$，只有两种形式：
		\begin{enumerate}
			\item $n = p^8$：只有一个质因子，指数加一等于 $9$。
			\item $n = p^2 \cdot q^2$（$p \ne q$）：两个质因子，$(2+1)(2+1) = 9$。
		\end{enumerate}

		\subsubsection{Legendre 公式（$n!$ 的质因子指数）}

		对质数 $p$，$n!$ 中 $p$ 的指数为
		$$v_p(n!) = \sum_{i = 1}^{\infty} \left\lfloor \frac{n}{p^i} \right\rfloor$$

		求和到 $p^i > n$ 自动截断，实际只跑 $O(\log_p n)$ 项。

		\textbf{常见特例}：
		\begin{itemize}
			\item \textbf{阶乘末尾零的个数}：$n!$ 末尾 $0$ 的个数 $= \min(v_2(n!), v_5(n!)) = v_5(n!)$（因为 $v_2 \ge v_5$），即 $\lfloor n/5 \rfloor + \lfloor n/25 \rfloor + \lfloor n/125 \rfloor + \cdots$。
			\item \textbf{组合数中某质因子的指数}：$v_p\!\left(\binom{n}{k}\right) = v_p(n!) - v_p(k!) - v_p((n-k)!)$。结合 \textbf{Kummer 定理}：$v_p\!\left(\binom{n}{k}\right)$ 等于 $p$ 进制下 $k + (n - k)$ 的进位次数。可用于快速判断 $\binom{n}{k}$ 的奇偶性（$p = 2$ 时进位次数 $= 0 \iff \binom{n}{k}$ 是奇数，进一步等价于 $k \mathbin{\&} (n - k) = 0$，即 $k$ 的二进制是 $n$ 的子集）。
		\end{itemize}

		\textbf{闭式}（Legendre 等价形式）：$v_p(n!) = \dfrac{n - s_p(n)}{p - 1}$，其中 $s_p(n)$ 是 $n$ 在 $p$ 进制下各位数字之和。

		\subsubsection{二次方程根与二阶线性递推}

		若 $\alpha, \beta$ 都是某个二次方程 $x^2 = px + q$ 的根，则 $S_n = k_1 \alpha^n + k_2 \beta^n$ 满足二阶线性递推：
		$$S_n = p \cdot S_{n-1} + q \cdot S_{n-2}$$

		经典例子：Fibonacci 数列 $F_n = \varphi^n + \psi^n$（$\varphi, \psi = \frac{1 \pm \sqrt{5}}{2}$），对应 $x^2 = x + 1$，递推为 $F_n = F_{n-1} + F_{n-2}$。

		\textbf{例}：$x_1, x_2 = 11 \pm 2\sqrt{30}$ 满足 $x^2 = 22x - 1$，因此 $S_n = x_1^n + x_2^n$ 的递推为
		$$S_n = 22 S_{n-1} - S_{n-2}$$

		\subsubsection{逆序对奇偶性与置换环}

		\textbf{交换一定改变逆序对数量的奇偶性}。长度为 $\mathrm{len}$ 的循环移动一格等价于 $\mathrm{len} - 1$ 次交换。

		例题：\href{https://codeforces.com/gym/105576/problem/D}{第 49 届 ICPC 区域赛沈阳站 D. Dot Product Game}

		\begin{minted}{cpp}
FOR(i, 1, N - 1) {
	char ch;
	ll l, r, d;
	cin >> ch >> l >> r >> d;
	ll len = r - l + 1;
	// 向左循环移动 d 次 = 交换 d*(len-1) 次
	if (d != 0)
		inva += (len - 1) * d;
	out();
}
		\end{minted}

		\textbf{快速判断排列逆序对的奇偶性}：将排列分解为置换环。一个长度为 $k$ 的环需要 $k - 1$ 次交换复位，因此逆序对数与 $n$ 减去环的数量奇偶性相同。设排列有 $c$ 个环，则逆序对数的奇偶性等于 $n - c$ 的奇偶性。比直接用树状数组求逆序对好写。

	\end{multicols*}

